\chapter{Cryptography — Schnorr Protocol}
%%% generated by the blueprint pipeline from TCSlib.Cryptography.SchnorrProtocol
\section{Overview}

This module formalises the Schnorr identification protocol over an arbitrary
commutative group $G$ with a generator $g$ of prime order $q$.  It establishes
the three core security properties: \emph{completeness} (honest transcripts
always pass verification), \emph{special soundness} (two accepting transcripts
with distinct challenges yield a witness extractor), and \emph{honest-verifier
zero knowledge} (simulator transcripts equal honest transcripts up to a
bijective reindexing of the randomness).

%%%%%%%%%%%%%%%%%%%%%%%%%%%%%%%%%%%%%%%%%%%%%%%%%%%%%%%%%%%%%%%%%%%%%%%%%%%%%%%
\section{Declarations}

\begin{definition}[Schnorr transcript]
\lean{Schnorr.Transcript}
\label{Schnorr.Transcript}
\leanok
A Schnorr transcript is a triple $(a, c, s) \in G \times \mathbb{Z}_q \times
\mathbb{Z}_q$, where $a$ is the prover's commitment (an element of the group
$G$), $c \in \mathbb{Z}_q$ is the verifier's challenge, and $s \in \mathbb{Z}_q$
is the prover's response.
\end{definition}

\begin{definition}[Prover commitment]
\lean{Schnorr.commit}
\label{Schnorr.commit}
\leanok
Given randomness $r \in \mathbb{Z}_q$ and a generator $g \in G$, the commitment
is $a := g^{r}$, where the exponent is taken via the canonical lift
$r.\mathrm{val} \in \mathbb{N}$.
\end{definition}

\begin{definition}[Prover response]
\lean{Schnorr.respond}
\label{Schnorr.respond}
\leanok
Given witness $w$, randomness $r$, and challenge $c$ in $\mathbb{Z}_q$, the
response is $s := r + c \cdot w \in \mathbb{Z}_q$.
\end{definition}

\begin{definition}[Verifier predicate]
\lean{Schnorr.Verify}
\label{Schnorr.Verify}
\leanok
The verifier accepts a transcript $(a, c, s)$ with respect to public key
$\mathit{pk} \in G$ when $g^s = a \cdot \mathit{pk}^{c}$, where the
exponents are taken via the canonical lifts to $\mathbb{N}$.
\end{definition}

\begin{definition}[Honest transcript]
\lean{Schnorr.honest}
\label{Schnorr.honest}
\leanok
\uses{Schnorr.Transcript, Schnorr.commit, Schnorr.respond}
The honest prover's full transcript for witness $w$, randomness $r$, and
challenge $c$ is the triple $(\mathrm{commit}(g, r),\; c,\; \mathrm{respond}(w, r, c))$,
i.e.\ $(g^r,\, c,\, r + c \cdot w)$.
\end{definition}

\begin{definition}[Simulator transcript]
\lean{Schnorr.simulate}
\label{Schnorr.simulate}
\leanok
\uses{Schnorr.Transcript}
The zero-knowledge simulator, given public key $\mathit{pk}$, challenge $c$,
and a fresh response $s$, outputs the transcript
$(g^s \cdot (\mathit{pk}^c)^{-1},\; c,\; s)$.
The commitment is chosen so that the transcript is accepting without using the
witness.
\end{definition}

\begin{definition}[Witness extractor]
\lean{Schnorr.extract}
\label{Schnorr.extract}
\leanok
Given two challenges $c_1, c_2$ and corresponding responses $s_1, s_2$ in
$\mathbb{Z}_q$, the extractor returns
\[
  \frac{s_1 - s_2}{c_1 - c_2} \;\in\; \mathbb{Z}_q.
\]
\end{definition}

\begin{theorem}[Completeness of the Schnorr identification protocol]
\lean{Schnorr.schnorr_completeness}
\label{Schnorr.schnorr_completeness}
\leanok
\uses{Schnorr.Verify, Schnorr.commit, Schnorr.respond}
\difficulty{4}
Let $q$ be a prime, let $G$ be a commutative group, and let $g \in G$ be an element of
order $q$. For any witness $w$, randomness $r$, and challenge $c$ in $\mathbb{Z}_q$, the
honest transcript — commitment $a = g^{r}$, challenge $c$, and response $s = r + c\,w$ —
is accepted by the verifier against the public key $\mathit{pk} = g^{w}$; that is,
\[
  g^{\,s} \;=\; a \cdot \mathit{pk}^{\,c},
\]
where each exponent is taken via the canonical lift of the corresponding element of
$\mathbb{Z}_q$ to $\mathbb{N}$.
\end{theorem}

\begin{theorem}[Special soundness of the Schnorr protocol]
\lean{Schnorr.schnorr_soundness}
\label{Schnorr.schnorr_soundness}
\leanok
\uses{Schnorr.Verify, Schnorr.extract}
\difficulty{4}
Special soundness. Let $q$ be prime, let $g$ be an element of order $q$ in a commutative
group $G$, and take the public key $\mathit{pk} = g^{w}$ for some witness $w \in
\bbz_q$. Suppose $a \in G$, and suppose $c_1, c_2 \in \bbz_q$ are distinct challenges
with responses $s_1, s_2 \in \bbz_q$ such that both transcripts $(a, c_1, s_1)$ and $(a,
c_2, s_2)$ are accepted by the verifier for $\mathit{pk}$. Then the extractor recovers
the witness:
\[
  \frac{s_1 - s_2}{c_1 - c_2} \;=\; w \quad\text{in } \bbz_q.
\]
\end{theorem}

\begin{definition}[Randomness reindexing]
\lean{Schnorr.reindex}
\label{Schnorr.reindex}
\leanok
For fixed $w, c \in \mathbb{Z}_q$, define the bijection
$\sigma_{w,c} : \mathbb{Z}_q \xrightarrow{\;\sim\;} \mathbb{Z}_q$ by
$r \mapsto r + c \cdot w$, with inverse $s \mapsto s - c \cdot w$.
\end{definition}

\begin{theorem}[Honest-verifier zero knowledge of the Schnorr protocol]
\lean{Schnorr.schnorr_hvzk}
\label{Schnorr.schnorr_hvzk}
\leanok
\uses{Schnorr.commit, Schnorr.honest, Schnorr.reindex, Schnorr.respond, Schnorr.schnorr_completeness, Schnorr.simulate}
\difficulty{3}
Let $G$ be a commutative group, let $q$ be prime, and let $g \in G$ be an element of
order $q$, taken as the public generator. Fix a witness $w \in \bbz_q$ and a challenge
$c \in \bbz_q$, and let $\mathit{pk} = g^{w}$ be the associated public key. Then, as
functions $\bbz_q \to G \times \bbz_q \times \bbz_q$ of the prover's randomness $r$, the
honest transcript map $r \mapsto \big(g^{r},\, c,\, r + c\cdot w\big)$ coincides with
the simulator map fed the reindexed randomness $\sigma_{w,c}(r) = r + c\cdot w$, namely
$r \mapsto \big(g^{\sigma_{w,c}(r)}\cdot(\mathit{pk}^{\,c})^{-1},\, c,\,
\sigma_{w,c}(r)\big)$; that is, these two maps are equal.
\end{theorem}
